\chapter{결론}

본 논문에서는 St. Petersburg 역설을 해결하기 위해 효용 함수를 사용하는 대신 시행 횟수-이익 확률 모델의 도입을 제안하였다. 이 모델은 이익 확률을 최대화하는 전략을 일반화한 의사결정 과정으로 이해할 수 있으며, 충분히 큰 시행 횟수에서는 기댓값 기반 의사결정과 동일한 결과를 도출하는 것을 확인하였다. 그러나 의사결정자가 무한한 반복 횟수를 고려하지 않는 현실에서는 이익 확률의 수렴 속도가 의사결정에 영향을 미치기 때문에, 성공 확률이 희박한 게임에 높은 참가 비용을 지불하지 않는 이유를 설명할 수 있었다. 결론적으로 기대효용이론의 한계를 보완할 수 있는 보조도구로 시행 횟수-이익 확률 모델을 사용함으로써, St. Petersburg 게임에서 높은 비용을 지불하지 않는 의사결정을 직관적으로 분석할 수 있었다.
